\documentclass[../../../include/open-logic-section]{subfiles}

\begin{document}

\olfileid{sth}{cardinals}{cardsasords}
\olsection{کاردینال‌ها به‌مثابه اعداد ترتیبی}

در واقع، نظریهٔ کاردینال‌های ما صرفاً (و بی‌پروا) از نظریهٔ اعداد
ترتیبی‌مان بهره خواهد گرفت. یعنی کاردینال‌ها را فقط به‌عنوان برخی
اعداد ترتیبیِ خاص تعریف خواهیم کرد. به‌طور مشخص، تعریف زیر را عرضه
می‌کنیم:

\begin{defn}\ollabel{defcardinalasordinal}
اگر بتوان $A$ را خوش‌ترتیب کرد، آن‌گاه $\card{A}$ کوچک‌ترین عدد
ترتیبیِ $\gamma$ است که $\cardeq{A}{\gamma}$. برای هر عدد ترتیبیِ
$\gamma$ می‌گوییم $\gamma$ یک \emph{کاردینال} است اگر و تنها اگر
$\gamma = \card{\gamma}$.
\end{defn}

همین حالا عبارت «می‌توان $A$ را خوش‌ترتیب کرد» را به کار بردیم.
چنان‌که تقریباً همیشه در ریاضیات رخ می‌دهد، تعبیر وجهی در اینجا فقط
شرحی غیررسمی از یک ادعای وجودی است: گفتنِ «می‌توان $A$ را خوش‌ترتیب
کرد» صرفاً یعنی «رابطه‌ای وجود دارد که $A$ را خوش‌ترتیب می‌کند».

اما \olref{defcardinalasordinal} مشکلی دارد. می‌خواهیم \emph{هر}
مجموعه اندازه‌ای داشته باشد؛ یعنی $\card{A}$ برای هر~$A$ وجود داشته
باشد. ولی تعریفی که همین حالا دادیم با یک شرطی آغاز می‌شود:
«\emph{اگر} بتوان $A$ را خوش‌ترتیب کرد\ldots». اگر مجموعه‌ای مانند
$A$ باشد که نتوان آن را خوش‌ترتیب کرد، تعریف ما به‌سادگی در تعریف
شیءِ~$\card{A}$ ناکام می‌ماند.

بنابراین، برای استفاده از \olref{defcardinalasordinal} به تضمینی
نیاز داریم که هر مجموعه را بتوان خوش‌ترتیب کرد. اما متأسفانه چنین
تضمینی در~$\ZF$ در دسترس نیست. پس اگر بخواهیم
\olref{defcardinalasordinal} را به کار ببریم، چاره‌ای جز افزودن اصل
موضوعی تازه، مانند اصل زیر، نداریم:
\begin{axiom}[خوش‌ترتیبی]
هر مجموعه را می‌توان خوش‌ترتیب کرد.
\end{axiom}
در \olref[choice][]{chap} بررسی خواهیم کرد که آیا اصل موضوع
خوش‌ترتیبی پذیرفتنی است یا نه. اما از این پس آزادانه آن را به کار
می‌بریم. با استفاده از آن، کاملاً سرراست می‌توان اثبات کرد که
کاردینال‌ها (به تعریفِ \olref{defcardinalasordinal}) وجود دارند و
به‌خوبی رفتار می‌کنند:

\begin{lem}\ollabel{lem:CardinalsExist}
برای هر مجموعهٔ $A$:
\begin{enumerate}
	\item\ollabel{cardaexists} $\card{A}$ وجود دارد و یکتاست؛
	\item\ollabel{cardaapprox} $\cardeq{\card{A}}{A}$؛
	\item\ollabel{cardaidem} $\card{A}$ یک کاردینال است؛ یعنی
	$\card{A} = \card{\card{A}}$.
\end{enumerate}
\end{lem}

\begin{proof}
$A$ را ثابت کنید. بنا بر اصل خوش‌ترتیبی، خوش‌ترتیبی‌ای مانند
$\tuple{A, R}$ وجود دارد. بنا بر
\olref[ordinals][ordtype]{thmOrdinalRepresentation}،
$\tuple{A, R}$ با عدد ترتیبیِ یکتایی مانند $\beta$ یکریخت است. پس
$\cardeq{A}{\beta}$. بنا بر استقرای ترامتناهی، کوچک‌ترین عدد ترتیبیِ
یکتایی مانند $\gamma$ وجود دارد که $\cardeq{A}{\gamma}$. بنابراین
$\card{A} = \gamma$؛ و این \olref{cardaexists} و
\olref{cardaapprox} را اثبات می‌کند. برای اثبات \olref{cardaidem}
توجه کنید که اگر $\delta \in \gamma$، آن‌گاه به سبب انتخاب
$\gamma$ داریم $\cardless{\delta}{A}$؛ و چون هم‌شماری رابطهٔ
هم‌ارزی است
(\olref[sfr][siz][equ]{equinumerosityisequi})،
$\cardless{\delta}{\gamma}$ نیز برقرار است. پس
$\gamma = \card{\gamma}$.
%برای برهان خلف، فرض کنید عدد ترتیبیِ $\delta \in \gamma$ای باشد که $\cardeq{\gamma}{\delta}$. آن‌گاه از هم‌ارزیِ هم‌شماری تناقضی با انتخاب $\gamma$ به‌عنوان کوچک‌ترین عدد ترتیبیِ هم‌شمار با $A$ حاصل می‌شود.
\end{proof}

نتیجهٔ بعدی اصل کانتور، و حتی بیش از آن، را تضمین می‌کند. (توجه
کنید که کاردینال‌ها ترتیب خود را از اعداد ترتیبی به ارث می‌برند؛
یعنی $\cardfont{a} < \cardfont{b}$ اگر و تنها اگر
$\cardfont{a} \in \cardfont{b}$. در صورت‌بندی این مطلب برای اشیایی
که می‌دانیم کاردینال‌اند از حروف فراکتور استفاده می‌کنیم. این
قرارداد نسبتاً رایج است. بدیل رایج دیگر استفاده از حروف یونانی است،
زیرا کاردینال‌ها عدد ترتیبی‌اند، اما آن حروف را از میانهٔ الفبا
برمی‌گزینند؛ برای نمونه: $\kappa, \lambda$.)
\begin{lem}\ollabel{lem:CardinalsBehaveRight}
برای هر دو مجموعهٔ $A$ و $B$:
\begin{align*}
	\cardeq{A}{B} &\text{ اگر و تنها اگر } \card{A} = \card{B}\\
	\cardle{A}{B} &\text{ اگر و تنها اگر } \card{A} \leq \card{B}\\
	\cardless{A}{B}&\text{ اگر و تنها اگر } \card{A} < \card{B}
\end{align*}
\end{lem}

\begin{proof}
جهت چپ‌به‌راستِ ادعای دوم را اثبات می‌کنیم (حالت‌های دیگر مشابه‌اند
و به‌عنوان تمرین واگذار می‌شوند). نمودار زیر را در نظر بگیرید:
\begin{center}
	\begin{tikzpicture}
	\node (nodea) {$A$};
	\node[right = 6em of nodea] (nodeb) {$B$};
	\node[below = 2em of nodea] (nodecarda) {$\card{A}$};
	\node[below = 2em of nodeb] (nodecardb) {$\card{B}$};
	\draw[->] (nodea)--(nodeb);
	\draw[<->] (nodea)--(nodecarda);
	\draw[<->] (nodeb)--(nodecardb);
	\draw[->, dashed] (nodecarda)--(nodecardb);
\end{tikzpicture}
\end{center}
پیکان‌های دوسر نشان‌دهندهٔ !!{bijection}هایی‌اند که وجودشان را
\olref{lem:CardinalsExist} تضمین می‌کند. با فرض
$\cardle{A}{B}$، !!a{injection} مانند $A\to B$ وجود دارد. اکنون با
دنبال‌کردن پیکان‌ها از $\card{A}$ به $A$، سپس به $B$ و بعد به
$\card{B}$، !!a{injection} از $\card{A}$ به $\card{B}$ به دست
می‌آوریم (پیکان خط‌چین).
\end{proof}\noindent
همچنین می‌توانیم با \olref{lem:CardinalsBehaveRight} قضیهٔ
شرودر--برنشتاین را دوباره اثبات کنیم. این قضیه می‌گوید اگر
$\cardle{A}{B}$ و $\cardle{B}{A}$، آن‌گاه $\cardeq{A}{B}$. آن را
به‌صورت \olref[sfr][siz][sb]{thm:schroder-bernstein} بیان کردیم، اما
نخست آن را ــ با قدری زحمت ــ در
\olref[sfr][infinite][card-sb]{sec} اثبات کردیم. اکنون ملاحظه کنید:

\begin{proof}[اثبات دوبارهٔ Schr\"oder--Bernstein]
اگر $\cardle{A}{B}$ و $\cardle{B}{A}$، آن‌گاه بنا بر
\olref{lem:CardinalsBehaveRight} داریم
$\card{A} \leq \card{B}$ و $\card{B} \leq \card{A}$. بنابراین، بنا
بر سه‌شقی و \olref{lem:CardinalsBehaveRight}،
$\card{A} = \card{B}$ و $\cardeq{A}{B}$.
\end{proof}
\noindent
هرچند این اثبات بسیار ساده است، به‌طور ضمنی هم به جایگزینی (برای
به‌دست‌آوردن
\olref[ordinals][ordtype]{thmOrdinalRepresentation}) و هم به اصل
خوش‌ترتیبی (برای تضمین \olref{lem:CardinalsBehaveRight}) وابسته
است. در مقابل، اثباتِ \olref[sfr][infinite][card-sb]{sec} بسیار
خودبسنده‌تر بود (در واقع، می‌توان آن را در~$\Zminus$ انجام داد).

\end{document}
